\section{Wnioski}

Celem niniejszej pracy magisterskiej było przeprowadzenie analizy porównawczej i ocena skuteczności metod obrony przed bezpośrednimi atakami promptowymi na duże modele językowe. W ramach badań zestawiono ze sobą trzy odmienne od siebie mechanizmy zabezpieczeń stosowanych na etapie wdrożenia modelu. Było to podejście oparte na inżynierii promptów Self-Reminder, wykorzystanie zewnętrznego klasyfikatora Llama-Guard oraz metoda oparta na inżynierii aktywacji. Ewaluację tych technik obrony przeprowadzono nie tylko wyłącznie na bazie skuteczności przeciwdziałania atakom, lecz skupiono sie również na aspektach takich jak analiza kompromisu między bezpieczeństwem a użytecznością modelu oraz wprowadzane przez nie opóźnienia.

Można dojść do wniosku, iż każda z zastosowanych metod mitygacji wymagała kompromisu między bezpieczeństwem, wydajnością systemu, użytecznością modelu oraz jakością generowanych treści. Metoda sterowania aktywacjami okazała się być najbardziej optymalnym rozwiązaniem pod wieloma względami. Technika ta zapewniła dla modelu Llama-3.1-8B-Instruct skuteczność obrony porównywalną do zewnetrznego klasyfikatora, bez jednoczesnego zapotrzebowania na dodatkowe zasoby pamięciowe, wytworzając znikome opóźnienie generacji pierwszego tokenu oraz zachowując ogólne zdolności modelu. Tym samym wzmocniono wbudowane w model mechanizmy bezpieczeństwa i nie naruszono w znaczącym stopniu jego standardowego działania. Na podstawie przeprowadzonych na modelu Mistral-7B-Instruct-v0.3 badań można stwierdzić, że sterowanie aktywacjami nie jest metodą uniwersalną dla każdej architektury i wymaga, by model docelowy posiadał podstawowo rozwinięty poziom bezpieczeństwa. W przeciwnym wypadku metoda ta nie zapewnia ochrony przed atakami.

Metoda mitygacji oparta na zewnętrznym modelu moderującym Llama-Guard zapewniła porównywalny poziom bezpieczeństwa dla dwóch odmiennych rodzin modeli, niezależnie od ich bazowej odporności. Generowała ona natomiast istotny narzut pamięciowy oraz obliczeniowy. Przełożyło się to dziesięciokrotny wzrost otrzymania pierwszego tokena w scenariuszu, w którym klasyfikowano prompty wejściowe. Przy zastosowaniu go do sprawdzania zapytań oraz także odpowiedzi modelu docelowego, wykluczona została możliwość strumieniowania tekstu do użytkownika oraz wytworzone zostało znaczne opóźnienie. Dla metody tej uzyskano również najwyższy odsetek nadmiernej odmowy na nieszkodliwe zapytania. Aspekty te mogą wykluczać użycie takiego podejścia w systemach czasu rzeczywistego, natomiast technika to może być akceptowalna w rozwiązaniach asynchronicznych. Zastosowanie Llama-Guard w systemach interaktywnych wymagałoby wykorzystania mocniejszych silników inferencyjnych, by zoptymalizować generowane opóźnienia, ewentualnie dokonywać analizy zapytania wejściowego partiami.

Technika Self-Reminder oparta na inżynierii promptów wytworzyła pomijalne opóźnienie porównywalne z metodą sterowania aktywacjami Natomiast okazała się nieskuteczna w obronie modelu przed atakami \gls{DAN} oraz \gls{PAIR} lecz skutecznie blokowała bezpośrednie szkodliwe instrukcje czy atak \gls{GCG}. Rozwiązanie to w największym stopniu wpłynęło na degradację ogólnych zdolności modelu, badanych na bazie testu MMLU.

Najskuteczniejszym narzędziem do omijania zabezpieczeń okazał się atak \gls{PAIR}. Uzyskał on wysoki poziom wskaźnika \gls{ASR} dla obydwu rozpatrywanych modeli. Choć wdrożone metody mitygacji znacząco zredukowały jego skuteczność względem bazowej podatności, nie zdołano zagwarantować pełnej ochrony przed tym typem ataku. Nawet w najlepszym wariancie obrony osiągnął on skuteczność równą 21\% dla modelu Llama oraz 33\% dla modelu Mistral. W przypadku braku zastosowania dodatkowych mechanizmów obronnych, wyjątkowo skuteczny okazał się atak \gls{DAN} opierający się na odgrywaniu ról i wprowadzeniu fikcyjnego kontekstu. Wymagał on nieporównywalnie mniej zasobów od innych wektorów ataku.

Przeprowadzona analiza dowodzi, iż żadna z metod mitygacji nie ma charakteru uniwersalnego. W związku z tym każda implementacja wymaga dostosowania do konkretnej architektury, docelowego przeznaczenia oraz środowiska wdrożenia. Wybór odpowiednich mechanizmów obronnych powinno stanowić wynik analizy, w której metryki bezpieczeństwa są maksymalizowane kosztem użyteczności i wydajności systemu. Ponadto, nieustanny rozwój nowych wektorów ataków podkreśla konieczność stałego doskonalenia zabezpieczeń oraz uzasadnia niezbędność prowadzenia dalszych badań w dziedzinie ochrony dużych modeli językowych.